% US-082 DE4 (Lower Saxony) transfer experiment --- the positive counterpart of the
% multi-region study. \input-able section: no \documentclass / \begin{document}.
% Prose in English. Every figure is a real EPIC 12 artefact; every number is anchored
% with a % src: comment to its source (report.json / engram-confirmed run).
% This section is \input by main.tex right after experiments_multiregion.tex.

\subsection{A Region Where Transfer Works: Lower Saxony (DE4)}
\label{sec:de4}

The multi-region study above is, by design, a catalogue of \emph{honest negatives}:
zero-shot fails across the Franco-Iberian and tropical gaps, and pooling regions widens
the taxonomy without rescuing a single fine class above production quality. A natural
question follows: \emph{is there a target region where the France-trained champion transfers
well, and what makes it different?} This subsection answers yes, and isolates the three
levers that decide it. The target is Lower Saxony, Germany (NUTS region DE4), chosen because
its EuroCrops~v2 coverage is \emph{complete} (not partial as in the Italian Tuscany pilot)
and its agriculture is dominated by large, phenologically separable winter cereals.

\paragraph{Building a PASTIS-homologous dataset for DE4.}
We materialise a dense PASTIS-homologous dataset for DE4 2023 from first principles, with
\textbf{100\% real data and zero placeholders}: per-patch Sentinel-2 L2A time series
(\texttt{DATA\_S2/S2\_*.npy}, 10 bands, $128\times128$) pulled from the Copernicus Data
Space Ecosystem (CDSE) Process API, dense crop masks (\texttt{ANNOTATIONS/TARGET\_*.npy})
burnt from the official EuroCrops~v2 polygons (JRC, HCAT4 taxonomy), and the 64-dimensional
AlphaEarth annual embedding (\texttt{SATELLITE\_EMBEDDING/V1/ANNUAL} v1.1, year 2023) pulled
via Google Earth Engine for every labelled parcel. The label space is the native German
HCAT crop set (37 classes), built by the same region-generalised builder used for Italy
(\texttt{--region-prefix de4}); every HCAT code is verified against the official JRC
crosswalk, none is invented.

\paragraph{Three levers, measured independently.}
Starting from the Tuscany pilot ceiling (Voting-3 macro F1~$\approx0.12$), we identified and
measured three independent levers, each on real out-of-fold data
(Table~\ref{tab:de4_levers}):
\begin{enumerate}[leftmargin=1.4em,itemsep=2pt]
  \item \textbf{Acquisition year 2023 vs.\ 2018}: on Tuscany, the XGBoost-on-AlphaEarth
    member rises from F1~$0.156$ to $0.225$ ($+44\%$) purely by moving the annual embedding
    to a more recent, better-covered year. % src: us-082-REPORTE-RESULTADOS.md Table 2
  \item \textbf{Temporal window of 14 months vs.\ 8}: PASTIS-R spans 13 months
    (Sep~2018--Oct~2019, including winter), whereas the Italian pilot covered only
    Mar--Oct (8 months), starving the winter-cereal phenology. Re-extracting with a
    \texttt{2022-11-01..2023-12-31} window (and a relaxed \texttt{max\_cloud}~$=50$) raises
    the usable date count from 24 to 56 on Tuscany and to $\sim$41 on DE4 --- close to the
    43 dates of PASTIS. % src: engram HALLAZGO VENTANA TEMPORAL US-082
  \item \textbf{Region DE4 vs.\ Tuscany}: the dominant lever. With large, clean parcels
    (median 10\,ha vs.\ 0.4\,ha) and full crop coverage, the Voting-3 macro F1 jumps from
    $0.119$ to $0.266$ ($2.2\times$), and the winter cereals --- wheat, maize, rye, rapeseed
    --- recover from F1~$0.05$--$0.20$ to $0.57$--$0.70$. % src: report.json de4_2023
\end{enumerate}

\begin{table}[t]
  \centering
  \caption{The three independent levers that lift transfer from the Tuscany ceiling to DE4,
  each measured on real out-of-fold data. The region change (Tuscany$\rightarrow$DE4) is the
  dominant factor: large separable winter-cereal fields close the gap that small fragmented
  Mediterranean parcels could not. XGBoost-on-AlphaEarth member unless noted.}
  \label{tab:de4_levers}
  \begin{tabular}{lcc}
    \toprule
    Configuration & XGB macro F1 & Accuracy \\
    \midrule
    Tuscany 2018 (8-month window, 24 dates)   & 0.156 & 0.37 \\
    Tuscany 2023 (14-month window, 56 dates)  & 0.225 & 0.45 \\
    DE4 2023 (14-month window, $\sim$41 dates) & \textbf{0.267} & \textbf{0.633} \\
    \bottomrule
  \end{tabular}
\end{table}

\paragraph{The champion Voting-3 on DE4.}
We re-run the full champion recipe on DE4 --- TSViT-pheno (Full-M capacity, warm-started from
the PASTIS champion) plus U-TAE plus XGBoost-on-AlphaEarth, aggregated to parcels and
combined by the three-member voting ensemble (\texttt{ItaliaParcelVotingEnsemble}). The
best dense member, TSViT-pheno-fullm, reaches fine-level F1~$0.333$ and accuracy~$0.682$ on
the held-out fold, and the Voting-3 scores macro F1~$0.266$ % src: report.json de4_2023
--- a $2.2\times$ improvement over the Tuscany Voting-3 ($0.119$). Crucially, the per-class
anatomy confirms the cardinality thesis of Section~\ref{sec:multiregion:perclass} from the
\emph{positive} side: where DE4 has large separable fields, the previously-collapsed cereals
rescue (maize~$0.703$, winter~rapeseed~$0.689$, winter~wheat~$0.674$, winter~rye~$0.573$),
yielding 8~classes at F1~$\ge0.60$ and 12 at $\ge0.40$, against a single class over $0.60$ on
Tuscany. % src: report.json de4_2023 per-class

\paragraph{Three taxonomic views (the ``three ways'').}
For an honest reading we evaluate the same trained ensemble under three label-space
conventions, decided a~priori: \emph{(A) conserving the native 37 HCAT classes}
(fine, macro~F1~$0.266$); \emph{(B) without conserving}, crosswalking the prediction onto the
PASTIS-style coarse space (coarse, macro~F1~$0.287$); and \emph{(C) the full end-to-end
procedure} (the trained vote itself, $0.266$). % src: report.json de4_2023 fine/coarse
The coarse view (B) is the most favourable, as expected when near-inseparable sibling cereals
are fused, but we report all three so the number is not cherry-picked by taxonomy.

\paragraph{A measured phenological branch for DE4.}
A subtle but important methodological point: the TSViT-pheno semantic branch aligns each
pixel against a per-class phenological prototype --- a mean-NDVI-curve description generated by
Gemini and encoded by a sentence transformer (after Wen et al.,
2025~\cite{wen2025phenology}). The prototypes are region-specific: a German winter-cereal
calendar, not the Mediterranean one. We generate the DE4 prototypes from the real per-class
NDVI curves of the DE4 patches, with the description grounded in the Lower-Saxony 2023
calendar; the curves are real and only the textual qualifier is regional, never fabricated.
This is the same discipline applied throughout: real signals, regional grounding, no
synthetic data.

\paragraph{Honest caveats.}
The DE4 Voting-3 is measured on a single held-out fold; the inattackable headline number
would require a 5-fold out-of-fold protocol ($5\times$ compute), although the per-class
pattern is decisive and would not change the verdict. Spatial coverage of the DE4 sample is
partial --- German parcels are large and dispersed, so the labelled patches cover a fraction
of the regional universe --- but the per-class loss is uniform across classes (no selection
bias), so the sample is representative. DE4 already wins decisively at this coverage; widening
it only improves the estimate. This is the positive bookend to the multi-region study: the
same recipe that fails zero-shot across distant gaps \emph{succeeds} when the target region
supplies large, phenologically separable fields and a complete, recent, winter-inclusive
temporal record.
